% GENERATED FILE. Edit canonical lessons, then run npm run generate:books.
% canonical-source-hash: fnv1a64:11fc8452b851022b
% canonical-lessons: TA-C519-oliperukki, TA-C519-olivanki, TA-C519-tolaipeci, TA-C519-accuppori, TA-C519-minnanu

\chapter{Sound and Print}
\label{ch:ta-sound-and-print}
\hlchaptermodality{519}

\begin{chapteropening}
\textbf{By the end of this chapter:} \emph{I can talk about speakers, microphones and printers.}

Complete the last lesson of chapter 519: I can talk about speakers, microphones and printers.
\end{chapteropening}

\section[\texorpdfstring{\ta{ஒலிபெருக்கி} (oliperukki)}{ஒலிபெருக்கி (oliperukki)}]{\ta{ஒலிபெருக்கி} (oliperukki) — a loudspeaker}

\label{lesson:TA-C519-oliperukki}



\begin{quote}
Before the new one: say the Tamil for a robot, then the Tamil for artificial.
\end{quote}

\subsection*{You'll want to know: \ta{ஒலிபெருக்கி}}
\textbf{\ta{ஒலிபெருக்கி}} — \emph{oliperukki} — \textquotedblleft{}a loudspeaker\textquotedblright{}.

\begin{quote}
\ta{கோயில்} \ta{திருவிழாவில்} \ta{ஒலிபெருக்கி} \ta{மிகவும்} \ta{சத்தமாக} \ta{இருந்தது}.
\end{quote}

\begin{quote}
\emph{kōyil tiruviḻāvil oliperukki mikavum cattamāka iruntatu.}
\end{quote}

\begin{quote}
At the temple festival the loudspeaker was very loud.
\end{quote}

\ta{ஒலி} \textquotedblleft{}sound\textquotedblright{} + \ta{பெருக்கி} \textquotedblleft{}that which multiplies\textquotedblright{}.

\begin{cousinweb}
\textbf{\ta{ஒலிபெருக்கி}} (\emph{oliperukki}) is a modern native compound: \ta{ஒலி} (\emph{oli}), \textquotedblleft{}sound\textquotedblright{}, and \ta{பெருக்கி} (\emph{perukki}), \textquotedblleft{}that which makes it great\textquotedblright{}, from \ta{பெரு} (\emph{peru}), \textquotedblleft{}great\textquotedblright{} (DEDR 4411).
\end{cousinweb}

\begin{grammarlens}[title={What you've built}]
One more word for describing what you lived and felt.
\end{grammarlens}

\subsection*{Your turn}
\begin{itemize}
\raggedright
  \item \emph{Say it:} \emph{oliperukki}
  \item \emph{Say it:} \emph{oliperukki}, once more
  \item \emph{Say it:} \emph{kōyil tiruviḻāvil oliperukki mikavum cattamāka iruntatu.} — At the temple festival the loudspeaker was very loud.
  \item \emph{Recall:} say the Tamil for a robot, then the Tamil for artificial, then say \emph{oliperukki} again
\end{itemize}

\begin{tcolorbox}[breakable,colback=teal!4,colframe=teal!35!black,title={Before you move on}]
What does \ta{ஒலிபெருக்கி} mean? (\textquotedblleft{}A loudspeaker\textquotedblright{}.) Say it once more.
\end{tcolorbox}
\section[\texorpdfstring{\ta{ஒலிவாங்கி} (olivāṅki)}{ஒலிவாங்கி (olivāṅki)}]{\ta{ஒலிவாங்கி} (olivāṅki) — a microphone}

\label{lesson:TA-C519-olivanki}



\begin{quote}
Before the new one: say the Tamil for artificial, then the Tamil for a loudspeaker.
\end{quote}

\subsection*{You'll want to know: \ta{ஒலிவாங்கி}}
\textbf{\ta{ஒலிவாங்கி}} — \emph{olivāṅki} — \textquotedblleft{}a microphone\textquotedblright{}.

\begin{quote}
\ta{பாடகர்} \ta{ஒலிவாங்கியைக்} \ta{கையில்} \ta{எடுத்தார்}.
\end{quote}

\begin{quote}
\emph{pāṭakar olivāṅkiyaik kaiyil eṭuttār.}
\end{quote}

\begin{quote}
The singer picked up the microphone.
\end{quote}

\ta{ஒலி} \textquotedblleft{}sound\textquotedblright{} + \ta{வாங்கி} \textquotedblleft{}that which takes in\textquotedblright{}, from \ta{வாங்கு} (\emph{vāṅku}).

\begin{cousinweb}
\textbf{\ta{ஒலிவாங்கி}} (\emph{olivāṅki}) is a modern native compound: \ta{ஒலி} (\emph{oli}), \textquotedblleft{}sound\textquotedblright{}, and \ta{வாங்கி} (\emph{vāṅki}), \textquotedblleft{}that which receives\textquotedblright{}, from \ta{வாங்கு}, \textquotedblleft{}to take, to get\textquotedblright{} (DEDR 5336).
\end{cousinweb}

\begin{grammarlens}[title={What you've built}]
One more word for describing what you lived and felt.
\end{grammarlens}

\subsection*{Your turn}
\begin{itemize}
\raggedright
  \item \emph{Say it:} \emph{olivāṅki}
  \item \emph{Say it:} \emph{olivāṅki}, once more
  \item \emph{Say it:} \emph{pāṭakar olivāṅkiyaik kaiyil eṭuttār.} — The singer picked up the microphone.
  \item \emph{Recall:} say the Tamil for artificial, then the Tamil for a loudspeaker, then say \emph{olivāṅki} again
\end{itemize}

\begin{tcolorbox}[breakable,colback=teal!4,colframe=teal!35!black,title={Before you move on}]
What does \ta{ஒலிவாங்கி} mean? (\textquotedblleft{}A microphone\textquotedblright{}.) Say it once more.
\end{tcolorbox}
\section[\texorpdfstring{\ta{தொலைபேசி} (tolaipēci)}{தொலைபேசி (tolaipēci)}]{\ta{தொலைபேசி} (tolaipēci) — a telephone; a landline}

\label{lesson:TA-C519-tolaipeci}



\begin{quote}
Before the new one: say the Tamil for a loudspeaker, then the Tamil for a microphone.
\end{quote}

\subsection*{You'll want to know: \ta{தொலைபேசி}}
\textbf{\ta{தொலைபேசி}} — \emph{tolaipēci} — \textquotedblleft{}a telephone; a landline\textquotedblright{}.

\begin{quote}
\ta{அலுவலகத்} \ta{தொலைபேசி} \ta{காலை} \ta{முதல்} \ta{அடிக்கிறது}.
\end{quote}

\begin{quote}
\emph{aluvalakat tolaipēci kālai mutal aṭikkiṟatu.}
\end{quote}

\begin{quote}
The office telephone has been ringing since morning.
\end{quote}

\ta{தொலை} \textquotedblleft{}far\textquotedblright{} + \ta{பேசி} \textquotedblleft{}that which speaks\textquotedblright{}. A mobile is a \ta{கைப்பேசி} (\emph{kaippēci}).

\begin{cousinweb}
\textbf{\ta{தொலைபேசி}} (\emph{tolaipēci}) is a modern native compound: \ta{தொலை} (\emph{tolai}), \textquotedblleft{}far, distant\textquotedblright{}, and \ta{பேசி} (\emph{pēci}), \textquotedblleft{}the speaker\textquotedblright{}, from \ta{பேசு} (\emph{pēcu}), \textquotedblleft{}to speak\textquotedblright{} (DEDR 4430).
\end{cousinweb}

\begin{grammarlens}[title={What you've built}]
One more word for describing what you lived and felt.
\end{grammarlens}

\subsection*{Your turn}
\begin{itemize}
\raggedright
  \item \emph{Say it:} \emph{tolaipēci}
  \item \emph{Say it:} \emph{tolaipēci}, once more
  \item \emph{Say it:} \emph{aluvalakat tolaipēci kālai mutal aṭikkiṟatu.} — The office telephone has been ringing since morning.
  \item \emph{Recall:} say the Tamil for a loudspeaker, then the Tamil for a microphone, then say \emph{tolaipēci} again
\end{itemize}

\begin{tcolorbox}[breakable,colback=teal!4,colframe=teal!35!black,title={Before you move on}]
What does \ta{தொலைபேசி} mean? (\textquotedblleft{}A telephone; a landline\textquotedblright{}.) Say it once more.
\end{tcolorbox}
\section[\texorpdfstring{\ta{அச்சுப்பொறி} (accuppoṟi)}{அச்சுப்பொறி (accuppoṟi)}]{\ta{அச்சுப்பொறி} (accuppoṟi) — a printer}

\label{lesson:TA-C519-accuppori}



\begin{quote}
Before the new one: say the Tamil for a microphone, then the Tamil for a telephone.
\end{quote}

\subsection*{You'll want to know: \ta{அச்சுப்பொறி}}
\textbf{\ta{அச்சுப்பொறி}} — \emph{accuppoṟi} — \textquotedblleft{}a printer\textquotedblright{}.

\begin{quote}
\ta{அச்சுப்பொறியில்} \ta{காகிதம்} \ta{தீர்ந்துவிட்டது}.
\end{quote}

\begin{quote}
\emph{accuppoṟiyil kākitam tīrntuviṭṭatu.}
\end{quote}

\begin{quote}
The printer has run out of paper.
\end{quote}

\ta{அச்சு} \textquotedblleft{}print, type\textquotedblright{} + \ta{பொறி} \textquotedblleft{}a machine\textquotedblright{}. The verb is \ta{அச்சிடு} (\emph{acciṭu}), to print.

\begin{cousinweb}
\textbf{\ta{அச்சுப்பொறி}} (\emph{accuppoṟi}) is a modern compound of \ta{அச்சு} (\emph{accu}), \textquotedblleft{}a mould, print\textquotedblright{}, and \ta{பொறி} (\emph{poṟi}), \textquotedblleft{}a machine\textquotedblright{}; the origin of both parts in these senses is not settled.
\end{cousinweb}

\begin{grammarlens}[title={What you've built}]
One more word for describing what you lived and felt.
\end{grammarlens}

\subsection*{Your turn}
\begin{itemize}
\raggedright
  \item \emph{Say it:} \emph{accuppoṟi}
  \item \emph{Say it:} \emph{accuppoṟi}, once more
  \item \emph{Say it:} \emph{accuppoṟiyil kākitam tīrntuviṭṭatu.} — The printer has run out of paper.
  \item \emph{Recall:} say the Tamil for a microphone, then the Tamil for a telephone, then say \emph{accuppoṟi} again
\end{itemize}

\begin{tcolorbox}[breakable,colback=teal!4,colframe=teal!35!black,title={Before you move on}]
What does \ta{அச்சுப்பொறி} mean? (\textquotedblleft{}A printer\textquotedblright{}.) Say it once more.
\end{tcolorbox}
\section[\texorpdfstring{\ta{மின்னணு} (miṉṉaṇu)}{மின்னணு (miṉṉaṇu)}]{\ta{மின்னணு} (miṉṉaṇu) — electronic}

\label{lesson:TA-C519-minnanu}



\begin{quote}
Before the new one: say the Tamil for a telephone, then the Tamil for a printer.
\end{quote}

\subsection*{You'll want to know: \ta{மின்னணு}}
\textbf{\ta{மின்னணு}} — \emph{miṉṉaṇu} — \textquotedblleft{}electronic\textquotedblright{}.

\begin{quote}
\ta{கடையில்} \ta{மின்னணுப்} \ta{பொருட்களுக்குத்} \ta{தள்ளுபடி} \ta{இருக்கிறது}.
\end{quote}

\begin{quote}
\emph{kaṭaiyil miṉṉaṇup poruṭkaḷukkut taḷḷupaṭi irukkiṟatu.}
\end{quote}

\begin{quote}
There is a discount on electronic goods in the shop.
\end{quote}

It stands before a noun and doubles a following stop: \ta{மின்னணுப்} \ta{பொருட்கள்}.

\begin{cousinweb}
\textbf{\ta{மின்னணு}} (\emph{miṉṉaṇu}) is a modern compound made to match English \emph{electronic}: the native \ta{மின்} (\emph{miṉ}), \textquotedblleft{}electric\textquotedblright{} (DEDR 4876), and \ta{அணு} (\emph{aṇu}), \textquotedblleft{}an atom\textquotedblright{}, from Sanskrit \emph{aṇu}, \textquotedblleft{}fine, minute\textquotedblright{} (CDIAL 193).
\end{cousinweb}

\begin{grammarlens}[title={What you've built}]
One more word for describing what you lived and felt.
\end{grammarlens}

\subsection*{Your turn}
\begin{itemize}
\raggedright
  \item \emph{Say it:} \emph{miṉṉaṇu}
  \item \emph{Say it:} \emph{miṉṉaṇu}, once more
  \item \emph{Say it:} \emph{kaṭaiyil miṉṉaṇup poruṭkaḷukkut taḷḷupaṭi irukkiṟatu.} — There is a discount on electronic goods in the shop.
  \item \emph{Recall:} say the Tamil for a telephone, then the Tamil for a printer, then say \emph{miṉṉaṇu} again
\end{itemize}

\begin{tcolorbox}[breakable,colback=teal!4,colframe=teal!35!black,title={Before you move on}]
What does \ta{மின்னணு} mean? (\textquotedblleft{}Electronic\textquotedblright{}.) Say it once more.
\end{tcolorbox}
